\section{Introduction}

\subsection{Industrial Background and Vibration-Based Fault Diagnosis}

Rotating machinery forms the backbone of critical infrastructure across manufacturing, energy production, aerospace, and transportation sectors worldwide. Essential components such as bearings, gears, shafts, and rotors are susceptible to progressive degradation due to mechanical wear, fatigue, overloads, and environmental stressors. Undetected faults can escalate into catastrophic failures, triggering unplanned downtime, equipment damage, safety hazards, and substantial economic losses. Industry analyses estimate global costs from such downtime exceeding \$50 billion annually in the United States alone, with predictive maintenance (PdM) strategies promising 20--50\% reductions in maintenance expenses \cite{cite1}\cite{cite22}. In heavy industries like oil \& gas and power generation, turbine and compressor failures contribute up to 30\% of outages, while automotive manufacturing reports bearing defects accounting for over 40\% of assembly line halts \cite{cite23}.

Vibration monitoring stands as the predominant non-intrusive diagnostic modality, leveraging accelerometers to capture oscillatory signals encoding structural dynamics and fault signatures. Healthy bearings produce characteristic broadband noise modulated by shaft rotational harmonics, whereas defective ones overlay impulsive excitations or modulated sidebands. Inner race defects generate high-amplitude impacts at the ball-pass frequency inner race (BPFI), while outer race faults yield stationary BPFO impacts \cite{cite2}. Early detection relied on demodulation techniques like envelope analysis and cepstral processing \cite{cite3}\cite{randall2011vibration}. Modern implementations deploy wireless sensor networks (WSNs) and edge computing for real-time surveillance, enabling condition-based maintenance \cite{lei2018review}. Nevertheless, signal interpretation demands specialized expertise to disentangle fault modulations from operational variabilities, highlighting the imperative for automated, semantically aware diagnostic systems \cite{cite4}\cite{cite24}.

\subsection{Data Scarcity and Domain Shift Challenges}

Contemporary fault diagnosis algorithms demand vast labeled datasets, yet industrial environments impose acute data scarcity. Controlled laboratory run-to-failure experiments are prohibitively expensive, generating limited faulty samples overshadowed by voluminous healthy baselines \cite{cite5}. Standard benchmarks like the Case Western Reserve University (CWRU) dataset \cite{casewestern2010} offer 10--12 fault severities under fixed conditions, but neglect real-world complexities such as multi-component faults or gradual degradation \cite{cite6}. The Paderborn University bearing dataset \cite{paderborn2014} introduces accelerated wear tests with diverse defects, yet remains constrained to specific operating regimes \cite{cite7}. Field data collection faces further hurdles: privacy regulations limit sharing proprietary sensor streams, sensor heterogeneity introduces inconsistencies, and rare severe faults skew class distributions toward imbalance ratios exceeding 1000:1 \cite{cite8}.

Deployment amplifies these issues through domain shifts. Laboratory setups employ idealized conditions, contrasting industrial sensors subjected to loose fixtures and thermal drifts \cite{cite25}. Operational variances—rpm fluctuations, torque cycles, lubrication changes—morph signal morphologies unpredictably, causing covariate shifts that degrade source-domain models by 30--60\% accuracy on targets \cite{cite9}\cite{cite26}.

\subsection{Limitations of Existing Methods}

Traditional signal processing laid the groundwork but exhibits inherent constraints. Fast Fourier Transform (FFT) resolves steady-state harmonics yet convolves transients into spectral smearing \cite{cite10}. Wavelet transforms provide time-frequency localization, but performance hinges on ad-hoc basis selection, rendering them brittle to noise \cite{cite11}\cite{dragomiretskiy2014emd}. Deep learning paradigms revolutionized the field: 1D CNNs on raw time-domain signals adeptly detect impulses \cite{cite12}\cite{wen2019convolutional}, while 2D CNNs applied to spectrograms exploit image-like invariances \cite{cite13}\cite{jia2016deep,verstraete2020cnn}. LSTMs and GRUs capture temporal evolutions but suffer gradient dilution over extended sequences \cite{cite14}. Recent CNN-Transformer hybrids leverage multi-head self-attention, achieving 98--99\% accuracies on intra-domain CWRU splits \cite{cite15}\cite{hoang2023vit}. Multimodal fusions with acoustic emissions or motor currents enhance resilience to single-sensor failures \cite{cite16}.

Despite advances, systemic shortcomings endure. Overfitting to surrogate lab distributions precipitates 20--50\% accuracy plunges across datasets like CWRU-to-Paderborn \cite{cite17}. Data augmentation techniques—GAN-generated synthetics or SMOTE oversampling—often inject domain-incongruent artifacts, eroding generalization \cite{cite18}. Metric-based few-shot learners (e.g., Prototypical Networks) falter on semantically divergent faults \cite{cite19}. Transfer learning and domain adaptation methods such as maximum mean discrepancy (MMD) minimization \cite{li2021mmd} and gradient reversal layers \cite{ganin2016domain} partially mitigate shifts but require target-domain statistics. Graph neural networks \cite{li2023gnn} and Transformer-based architectures \cite{vaswani2017attention}\cite{li2022transformer} improve intra-domain accuracy yet remain brittle to unseen fault categories. Critically, unimodal approaches forfeit exploitable fault ontologies and causal semantics, impeding interpretability.

Emerging vision-language models (VLMs) like CLIP \cite{radford2021clip}, BLIP-2, and LLaVA \cite{liu2023llava} demonstrate prowess in aligning natural images with textual descriptions via contrastive pretraining. Audio-language models such as AudioCLIP \cite{guzhov2022audioclip} and CLAP \cite{elizalde2023clap} extend this paradigm to non-visual signals, inspiring vibration-semantic alignment. However, their applicability to vibration analysis is curtailed by domain gaps: spectrograms deviate from natural imagery with engineered axes and absence of semantic priors in pretraining corpora \cite{cite27}\cite{cite28}. Recent works on VLM-based fault diagnosis \cite{liu2024faultgpt}\cite{zhang2024mmfault,li2024llm4fd} highlight both the promise and limitations of direct VLM application to industrial signals. MLLMs struggle with quantitative signal reasoning, exhibiting hallucination rates up to 40\% on technical prompts without modality-specific finetuning \cite{cite29}. These limitations underscore the need for tailored alignment strategies bridging vibrations to semantic fault narratives.

\subsection{VibSemAlign Overview and Contributions}

VLMs furnish a compelling paradigm for vibration diagnostics by projecting spectrograms—derived via short-time Fourier transform (STFT) as 2D heatmaps—into unified embedding spaces alongside textual fault descriptors \cite{cite20}\cite{cite21}. VibSemAlign harnesses this potential through targeted finetuning of a pretrained VLM backbone with a vibration-semantic contrastive loss, augmented by cross-attention for refined multimodal fusion. Model-agnostic meta-learning (MAML) further endows rapid adaptation to novel domains using only 5--10 target shots. Figure~\ref{fig:1} illustrates the comprehensive pipeline spanning raw vibration acquisition to probabilistic fault inference.

This paper delineates five principal contributions:
\begin{enumerate}
\item Introducing VibSemAlign, a pioneering MLLM framework enabling zero-shot diagnosis with 96.5\% accuracy on CWRU (12\% absolute gain over top CNN baselines) and 92.3\% on Paderborn (18\% improvement) \cite{cite30}\cite{cite31}.
\item Devising a domain-specialized contrastive loss surpassing standard NT-Xent by 8\% in embedding alignment, enforcing invariances to speed/load perturbations \cite{cite32}\cite{cite33}.
\item Seamlessly integrating MAML, reducing few-shot adaptation shots by 50--70\% compared to vanilla finetuning \cite{cite34}\cite{cite35}.
\item Comprehensive empirical scrutiny on benchmark datasets, including t-SNE projections evidencing tight semantic clustering (Silhouette score 0.78 vs. 0.52 for unimodal CNNs) and ablation isolating each module's impact \cite{cite36}\cite{cite37}.
\item Open-sourcing models and code, alongside standardized evaluation protocols to propel reproducible advances in semantic signal processing \cite{cite38}\cite{cite39}.
\end{enumerate}

Section II reviews related literature. Section III expounds the proposed methodology. Section IV presents experimental results, with conclusions in Section V.
